\section{Spazio di Hilbert}
Definiamo lo \textbf{spazio di Hilbert} come uno spazio di vettori $\mathcal{H}$, definito su $\mathbb{R}$ o $\mathbb{C}$, il quale ha un prodotto
interno: \[
    \left<.,. \right> : \mathcal{H}\times \mathcal{H} \to \mathbb{C}
\]

\noindent Il prodotto, supposto che sia lineare nel suo primo argomento, soddisfa alcune proprietà: \begin{itemize}
    \item \textbf{Linearità}: Per $z,z',\omega\in \mathcal{H}$ e $\alpha, \beta\in \mathbb{C}$, vale $\left<\alpha z + \beta z', \omega\right> =
    \alpha\left<z, \omega \right> + \beta\left<z', \omega \right>$
    \item \textbf{Simmetria coniugata}: $\left<z, z' \right> = \overline{\left<z', z \right>}$
    \item \textbf{Positività}: $\left<z, z \right> \geq 0$
    \item \textbf{Definitività positiva}: $\left<z, z \right> = 0 \iff z = 0$
\end{itemize}

\noindent Il prodotto interno induce anche una norma $||z|| = \sqrt{\left< z, z\right>}$, rispetto alla quale lo spazio Hilbertiano è completo. Questa
sua caratteristica è anche ciò che lo distingue da altri spazi semplici di prodotti interni. Inoltre, in quanto completo, ogni sequenza di Cauchy $z_n$
di elementi di $\mathcal{H}$, definita come: \[
    ||z_n - z_m|| \to 0\quad \text{con n, m} \to +\infty
\]

\noindent convergerà, secondo la norma, a un elemento appartenente allo spazio $\mathcal{H}$, il quale, per segnali unidimensionali, è considerato come
$\mathcal{H} = L^2(\mathbb{R};\mathbb{C})$, i cui elementi sono funzioni misurabili che associano un valore dei reali a uno nei complessi e soddisfano:
\[
    \int_{-\infty}^{+\infty} |z(t)|^2 dt < +\infty
\]

\noindent Concludendo che un segnale complesso a energia finita appartiene allo spazio dei segnali $z\in L^2(\mathbb{R};\mathbb{C})$, la cui forma
vettoriale è data da: \[
    (\alpha z + \beta z')(t) = \alpha z(t) + \beta z'(t)
\]

\noindent Dove $z, z'$ sono i vettori dei segnali, mentre $z(t), z'(t)$ sono numeri complessi che indicano il valore assunto dal vettore segnale al
tempo $t$.\par
Nel contesto della teoria dei segnali, il \textbf{prodotto interno} dei segnali può essere interpretato come l'\textbf{energia} del segnale. Questo
perché, a partire dalla forma del prodotto interno standard nello spazio hilbertiano, otteniamo la norma: \[
    \left<z, z' \right> = \int_{-\infty}^{+\infty} z(t)\overline{z'(t)} dt \implies ||z|| = \sqrt{\left<z, z \right>}
\]

\noindent Di conseguenza, se rimuoviamo la radice elevando la scrittura al quadrato, otteniamo: \[
    ||z||^2 = \left<z, z \right> = \int_{-\infty}^{+\infty} z(t)\overline{z(t)} dt
\]

\noindent E siccome vale $z(t)\overline{z(t)} = |z(t)|^2$, possiamo concludere la tesi, dimostrando che il prodotto interno è interpretabile come
l'energia del segnale, la cui simmetria è garantita dal complesso coniugato.\par
Negli spazi di Hilbert, quindi anche negli Euclidei, un vettore può essere scritto come la combinazione dei suoi elementi lungo gli assi cartesiani,
ottenuti attraverso proiezioni: \[
    x = x_1e_1 + x_2e_2 + ... \implies x_k = \left< x, e_k\right>
\]

\noindent Dove l'insieme $e_k$ è una base ortonormale per lo spazio, detta \textbf{base canonica}. Similmente, nello spazio hilbertiano abbiamo: \[
    z(t) = \sum_{k}c_k\phi_k(t)
\]

\noindent Dove il segnale generico $z(t)\in \mathcal{H}$ è definito come combinazione lineare tra una base ortonormale $\phi_k$ e i coefficienti $c_k$.
Diciamo inoltre, che la base è ortonormale se: \[
    \left<\phi_k, \phi_m \right> = \delta_{km}\quad \text{con } \delta_{km} = \begin{cases}
        1, k=m\\
        0, k\neq m
    \end{cases}
\]

\noindent Ogni funzione base ha quindi una norma unitaria $||\phi_k|| = 1$ e i coefficienti sono dati da $c_k = \left<z, \phi_k \right>$.

%

\section{Fourier}
Introdotto lo spazio di lavoro, possiamo definire innanzitutto la \textbf{base di Fourier} $\{\phi_k(t)\}$, la quale è composta dalle funzioni di base:
\[
    \phi_k(t) = \frac{1}{\sqrt{T}}e^{j2\pi kt/T}\quad \text{con }k\in\mathbb{Z}
\]

\noindent Per la quale ogni funzione di base è parametrizzata da $\omega_k = \frac{2\pi k}{T}$. Si può inoltre notare come appartenga alla classe di
funzioni sinusoidali complesse $z_k(t) = Ae^{j(\omega_kt+\phi_0)}$, dove: \begin{itemize}
    \item \textbf{Ampiezza}: $A = \frac{1}{\sqrt{T}}$
    \item \textbf{Frequenza angolare}: $\omega_k = \frac{2\pi k}{T}$
    \item \textbf{Fase iniziale}: $\phi_0 = 0$
    \item \textbf{Parte reale}: $Re\{\phi_k(t)\} = \frac{1}{\sqrt{T}}cos\left(\frac{2\pi kt}{T}\right)$
    \item \textbf{Parte immaginaria}: $Im\{\phi_k(t)\} = \frac{1}{\sqrt{T}}sin\left(\frac{2\pi kt}{T}\right)$
    \item \textbf{Modulo}: $|\phi_k(t)| = \frac{1}{\sqrt{T}}$
    \item \textbf{Fase}: $arg(\phi_k(t)) = \omega_kt = \frac{2\pi kt}{T}$
\end{itemize}

\noindent $T$, come già menzionato, è detto periodo fondamentale, per il quale sono definite anche le frequenze fondamentali che ritornano le frequenze
della base di Fourier: \[
    \omega_k = \frac{2\pi k}{T} = k\omega_0\quad f_k = \frac{k}{T} = kf_0
\]

\noindent Nominiamo, in questi termini $k=1$ la \textbf{frequenza fondamentale}, annotata con $\omega_1 = 1\omega_0 = \omega_0$, e con valori di $k$
maggiori di $1$ diremo che sono la seconda, terza, quarta \textbf{armonica}, e così via. Invece, per quanto riguarda valori negativi, il segnale è
rappresentato da una componente dalla frequenza negativa $\omega_{-k} = -\omega_k$. Dunque possiamo dire che $\phi_k(t)$ è la k-esima funzione di base
complessa armonica.\par
La base è interpretabile geometricamente per ogni valore fissato $t$; il numero complesso dato dalla base è un punto sul cerchio $|z| =
\frac{1}{\sqrt{T}}$, e al variare di $t$, ruota il punto intorno all'origine mantenendo un raggio costante uguale a $\frac{1}{\sqrt{T}}$ con posizione
angolare data da $\frac{2\pi kt}{T}$. In particolare: \begin{itemize}
    \item Con $k>0$ si ha una rotazione antioraria.
    \item Con $k<0$ si ha una rotazione oraria.
    \item Con $k=0$ non si ha alcuna rotazione.
\end{itemize}

\noindent Infine, le funzioni di \textbf{base di Fourier} formano una \textbf{base ortonormale} dello spazio $L^2([0,T], \mathbb{C})$, infatti: \[
    \left<\phi_k, \phi_m \right> = \int_{0}^{T} \phi_k(t)\overline{\phi_m(t)} dt = \delta_{km}
\]

\noindent Il risultato è detto \textbf{delta di Kronecker}, e rispetta le proprietà dell'ortonormalità; mostrando contemporaneamente che le funzioni di
base hanno norma unitaria $||\phi_k|| = 1$ e sono mutuamente ortogonali, dunque l'insieme $\{\phi_k\}$ è una base ortonormale e per ogni
$z(t)\in L^2([0,T], \mathbb{C})$ possiamo scrivere: \[
    z(t) = \sum_{k\in \mathbb{Z}} c_k\phi_k(t)
\]

\noindent Per definizione, il delta di Kronecker può assumere solo i valori discreti $0,1$, e per questa ragione ne si fa uso nei segnali digitali. Per
quanto riguarda i valori continui, invece, si utilizza il \textbf{delta di Dirac}, la cui funzione è la medesima, ma è capace di assumere anche valori
non discreti, rendendolo utile per lo studio dei segnali nel dominio del tempo.\par
Consideriamo ora un segnale periodico $z(t)\in L^2([0,T], \mathbb{C})$, è possibile espandere il segnale nella forma discussa pocanzi: \[
    z(t) = \sum_{k = -\infty}^{+\infty} c_k\phi_k(t)
\]

\noindent Dove i coefficienti sono le proiezioni di $x$ sulle funzioni di base di Fourier: $c_k = \left<x, \phi_k \right>$. Sostituendo la definizione
del prodotto interno otteniamo la \textbf{serie di Fourier}: \[
    c_k = \frac{1}{\sqrt{T}} \int_0^T x(t)e^{-j\omega_k t} dt
\]

\noindent Ogni coefficiente $c_k$ indica quanto la frequenza $\omega_k = \frac{2\pi k}{T}$ è presente all'interno del segnale; in altre parole
rappresenta la proiezione del segnale sulla sinusoide di frequenza $\omega_k$\par
Grazie alla serie di Fourier è possibile fare delle inferenze in merito alle componenti della scrittura. In particolare, se consideriamo la componente
k-esima della serie $z_k(t) = c_k\phi_k(t)$ e scriviamo il coefficiente complesso di Fourier in forma polare, otteniamo: \[
    c_k = |c_k|e^{jarg(c_k)}
\]

\noindent E da questo ricaviamo: \begin{equation}
    \begin{split}
        z_k(t) &= c_k\phi_k(t)\\
            &= |c_k|e^{jarg(c_k)}\cdot \frac{1}{\sqrt{T}}e^{j\omega_k t}\\
            &= \frac{|c_k|}{\sqrt{T}}e^{j(\omega_kt + arg(c_k))}
    \end{split}
\end{equation}

\noindent Dunque $z_k(t)$ è una sinusoide complessa della forma $A_ke^{j(\omega_kt + \phi_{0,k})}$, dove: \begin{itemize}
    \item \textbf{Ampiezza}: $A_k = |z_k(t)| = |c_k||\phi_k(t)| = \frac{|c_k|}{\sqrt{T}}$
    \item \textbf{Fase}: $arg(z_k(t)) = \omega_kt + arg(c_k)$
    \item \textbf{Fase iniziale}: $\phi_{0,k} = arg(c_k)$
\end{itemize}

\noindent Inoltre, sapendo che $c_k$ è un numero complesso; $|c_k|$ controlla la scala del componente e $\omega_kt$ fornisce la fase dipendente dal
tempo. In questi termini, la funzione $c_k\phi_k(t)$ è detta componente armonica k-esima del segnale.\par
Quando parliamo di segnali non periodici, il periodo fondamentale $T$ tende idealmente a infinito, avvicinando sempre di più le frequenze discrete
$\omega_k = \frac{2\pi k}{T}$, eventualmente facendo loro formare un insieme continuo. In questi termini, la sommatoria della serie di Fourier diventa
un integrale della forma: \[
    x(t) = \frac{1}{\sqrt{2\pi}}\int_{-\infty}^{+\infty} X(\omega)e^{i\omega t}d\omega
\]

\noindent Dal quale possiamo definire la \textbf{trasformata di Fourier}, un operatore che prende una funzione in input e ne ritorna il prodotto
scalare tra di essa e la funzione sinusoide complessa di frequenza $\omega$: \[
    X(\omega) = \frac{1}{\sqrt{2\pi}}\int_{-\infty}^{+\infty} x(t)e^{-i\omega t} dt
\]

\noindent In particolare, la funzione $\phi_\omega(t) = \frac{1}{\sqrt{2\pi}}e^{i\omega t}$ rappresenta le funzioni di base associate alla frequenza
$\omega$. Formalmente si scrive: \[
    X(\omega) = \left<x, \phi_\omega \right>
\]

\noindent E indica che la $X(\omega)$ è la proiezione del segnale sulla sinusoide complessa di frequenza $\omega$. Attenzione, tuttavia, che le funzioni
$e^{i\omega t}$ non appartengono allo spazio $L^2(\mathbb{R}, \mathbb{C})$, dunque, nel caso non periodico, non sono funzioni di base ordinarie dello
spazio di Hilbert, bensì quello che fanno è formare una base generalizzata; per questa ragione è più corretto dire che la trasformata di Fourier è la
\textbf{rappresentazione} di un segnale rispetto a una famiglia continua di funzioni di base generalizzata, come quella delle sinusoidi complesse.\par
In ogni caso, facendo uso della trasformata, è possibile passare dal dominio del tempo, i cui segnali sono in forma $x(t)$ a quello delle frequenze
$X(\omega)$, dove ogni valore è una proiezione: $X(\omega) = \left<x, \phi_\omega \right>$. Per ricostruire il segnale, invece, è necessario possedere i
valori di ampiezza $|X(\omega)|$ e fase $arg(X(\omega))$, per poi usare l'\textbf{antitrasformata}: \[
    x(t) = \int_{-\infty}^{+\infty} X(\omega)\phi_\omega(t) d\omega = \sum_k c_k\phi_k(t)
\]

\noindent Inoltre, siccome la trasformata di Fourier è un operatore unitario, è preservato il prodotto interno e, di conseguenza, la norma: \[
    \int_{-\infty}^{+\infty} |x(t)|^2 dt = \int_{-\infty}^{+\infty} |X(\omega)|^2 d\omega
\]

\noindent Questa scrittura è detta \textbf{relazione di Parseval} sotto convenzione simmetrica, e la quantità $|X(\omega)|^2$ è interpretabile come
densità di energia rispetto alla frequenza, rendendo calcolabile l'energia totale del segnale anche nel dominio delle frequenze. Se si vogliono
considerare solo le frequenze contenute in una certa \textbf{banda} $B$, è possibile moltiplicare lo spettro per una funzione indicatore: \[
    X_B(\omega) = X(\omega)l_B(\omega)
\]

\noindent Che fornirà il segnale filtrato: \[
    x_B(t) = \frac{1}{\sqrt{2\pi}}\int_B X(\omega)e^{i\omega t}d\omega
\]

\noindent Il quale è una proiezione sul subspazio generato dalle sinusoidi le cui frequenze si trovano all'interno della banda. Un filtro ideale è
considerabile come un operatore di proiezione sullo spazio delle frequenze.\par
Conoscere il contenuto frequenziale dei segnali grazie a Fourier è fondamentale per poter processarli; consente infatti di interpretare i segnali grazie
a una maggiore facilità di manipolazione, data dai valori discreti. Nel dominio del tempo, infatti, non è possibile avere un segnale "puro", perché di
norma è disturbato da altri. Il dominio delle frequenze risolve questo problema fornendo un ambiente privo di contaminazioni.\par
Esiste inoltre una relazione tra la forma del segnale nel dominio dello spazio e delle frequenze. Definiamo \textbf{larghezza spaziale}\footnote{Da non
confondersi con la frequenza spaziale, che misura quanto velocemente un segnale oscilla per ogni unità di distanza} $L$ di un segnale quella quantità
che indica quanto la funzione si estende lungo la coordinata spaziale del tempo $t$ ed è compresa nell'intervallo $[-L/2, L/2]$. La \textbf{larghezza
spettrale} (o frequenziale) è il valore analogo all'interno del dominio delle frequenze.\par
Queste due misure sono inversamente proporzionali: un segnale con una larghezza spaziale stretta nel dominio del tempo avrà un'ampia larghezza spettrale
e viceversa. Esistono anche dei casi particolari: \begin{itemize}
    \item Una funzione costante nello spazio diventa un impulso nel dominio delle frequenze.
    \item Una funzione rettangolare di lunghezza $L$ diventa una funzione $sinc$, la quale ha il suo primo zero in $\omega = \pm 1/L$.
    \item Una gaussiana con un certo $\sigma$ nel tempo ne diventa un'altra con $\sigma_\omega = 1/\sigma_t$.
\end{itemize}

%

\section{Convoluzione}
La \textbf{convoluzione} è un'operazione importantissima all'interno del campo della teoria dei segnali; descrive come un segnale di input $f$ è
modificato da un sistema la cui risposta è $g$. Siano quindi queste due funzioni appartenenti al dominio dei complessi, la definizione è: \[
    (f*g)(t) = \int_{-\infty}^{+\infty} f(\tau)g(t-\tau) d\tau
\]

\noindent Graficamente ha luogo tenendo fermo uno dei due segnali, di norma $f$, e si inverte $g$ per poi farlo scorrere sul primo, formando un'area
in base al tempo $t$ considerato. Geometricamente, dunque, misura quanto le due funzioni si sovrappongono quando una scorre sopra l'altra. Detiene
inoltre alcune proprietà: \begin{itemize}
    \item \textbf{Commutativa}: $f*g = g*f$
    \item \textbf{Associativa}: $(f*g)*h = f*(g*h)$
    \item \textbf{Distributiva rispetto all'addizione}: $f*(g+h) = f*g + f*h$
    \item \textbf{Linearità rispetto allo scalare}: $(af)*g = a(f*g)$
\end{itemize}

\noindent Introduciamo ora il \textbf{teorema della convoluzione}, il quale la mette in relazione con la trasformata di Fourier; afferma che la
convoluzione di due segnali nel dominio del tempo è equivalente al loro prodotto nel dominio delle frequenze: \[
    \widehat{f*g} = \sqrt{2\pi}\cdot F(\omega)G(\omega)
\]

\noindent Per dimostrarlo, applichiamo la trasformata di Fourier: \begin{equation}
    \begin{split}
        \widehat{f*g}(\omega) &= \frac{1}{\sqrt{2\pi}} \int_{-\infty}^{+\infty} (f*g)(t)\cdot e^{-i\omega t} dt\\
		&= \frac{1}{\sqrt{2\pi}} \int_{-\infty}^{+\infty} \int_{-\infty}^{+\infty} f(\tau)g(t-\tau)\cdot e^{-i\omega t} d\tau dt
    \end{split}
\end{equation}

\noindent Cambiamo ora qualche variabile negli integrali; diciamo $s = t-\tau, t = st\tau$, ottenendo: \[
    e^{-i\omega t} = e^{-i\omega(s+\tau)} = e^{-i\omega s}e^{-i\omega\tau}
\]

\noindent Gli integrali ora si possono separare e per la definizione della trasformata, ognuno è uguale a $\sqrt{2\pi}$ volte la sua trasformata
normalizzata: \begin{equation}
    \begin{split}
        \widehat{f*g}(\omega) &= \frac{1}{\sqrt{2\pi}} \left(\int_{-\infty}^{+\infty} f(\tau)e^{-i\omega\tau}d\tau\right) \left(\int_{-\infty}^{+\infty}
                g(s)e^{-i\omega s}ds\right)\\
            &= \frac{1}{\sqrt{2\pi}}\cdot (\sqrt{2\pi}F(\omega)) (\sqrt{2\pi}G(\omega))\\
            &= \sqrt{2\pi}\cdot F(\omega)G(\omega)
    \end{split}
\end{equation}

\noindent A sua volta, è possibile anche applicare la convoluzione nel dominio delle frequenze, la quale corrisponderà al prodotto dei segnali nel
dominio del tempo: \[
    (F*G)(\omega) = \int_{-\infty}^{+\infty}F(\nu)G(\omega-\nu)d\nu
\]

\noindent Considerando poi l'antitrasformata di Fourier, otteniamo: \[
    \mathcal{F}^{-1}\{F*G\}(t) = \sqrt{2\pi}\cdot f(t)g(t) \implies \widehat{f\cdot g}(\omega) = \frac{1}{\sqrt{2\pi}} (F*G)(\omega)
\]

\noindent La dimostrazione inizia prendendo la trasformata del prodotto e la sostituzione della trasformata inversa per $g(t)$: \begin{equation}
    \begin{split}
        \widehat{f\cdot g}(\omega) &= \frac{1}{\sqrt{2\pi}} \int_{-\infty}^{+\infty} f(t)g(t) e^{-i\omega t} dt\\
            &= \frac{1}{\sqrt{2\pi}} \int_{-\infty}^{+\infty} f(t) \left[\int_{-\infty}^{+\infty} G(\nu)e^{i\nu t} d\nu\right]e^{-i\omega t} dt
    \end{split}
\end{equation}

\noindent Si procede invertendo l'ordine di integrazione, ottenendo un integrale di Fourier non normalizzato di $f$ nelle quadre. Siccome la formula
della trasformata include la frazione, la si aggiunge nella trasformazione, ottenendo: \begin{equation}
    \begin{split}
        \widehat{f\cdot g}(\omega) &= \frac{1}{\sqrt{2\pi}} \int_{-\infty}^{+\infty} G(\nu) \left[\int_{-\infty}^{+\infty} f(t)e^{-i(\omega-\nu) t} dt
                \right]d\nu\\
            &= \frac{1}{\sqrt{2\pi}} \int_{-\infty}^{+\infty} G(\nu)\cdot F(\omega-\nu) d\nu = \frac{1}{\sqrt{2\pi}} (F*G)(\omega)
    \end{split}
\end{equation}

%

\section{Sistemi e filtri}
Definiamo \textbf{sistema} un'entità capace di ricevere in input un segnale $x(t)$ e produrre un output $y(t)$, chiamata risposta rispetto all'input.
Ogni sistema ha al suo interno un elemento $T$ che trasforma il segnale da input ad output e può essere un dispositivo fisico, un'operazione matematica,
un filtro o molto altro ancora. Detengono alcune proprietà fondamentali:

\paragraph{Linearità:} Detta anche \textbf{superposition}, afferma che sommare o scalare l'input provoca la somma o la scalatura anche nell'output. \[
	T\{ax_i + bx_2\} \equiv aT\{x_i\} + bT\{x_2\}; \text{Con $a,b$ scalari}
\]

\noindent L'utilità di questa proprietà è fondata sul fatto che input e output possono appartenere a domini diversi; di conseguenza, se si modifica uno
dei due segnali, così farà anche l'altro senza ulteriori operazioni.

\paragraph{Tempo-invarianza:} Aggiungendo un ritardo (time-shift) al segnale di input, lo si provocherà anche in quello di output: \[
	T\{x(t-t_0)\} = y(t-t_0)
\]

\noindent Chiariamo che questa proprietà non è universale; altri sistemi possono presentare una tempo-varianza, dove in base si troverà nel tempo il
segnale di input, quello di output potrebbe essere amplificato o modificato. Per esempio $x(t-t_0) \implies tx(t-t_0)$.

\paragraph{Causalità:} Un sistema si dice causale se il suo output al tempo $t$ dipende solo dal segnale di input fino al tempo $t$ e mai dopo.
Fondamentale per poter permettere una trasformazione in tempo reale. \begin{equation}
    \begin{aligned}
        y(t) &= 2x(t) & \text{Sistema causale}\\
        y(t) &= 2x(t-1) & \text{Sistema causale}\\
        y(t) &= x(t+1) & \text{Sistema non causale}
    \end{aligned}
\end{equation}

\noindent L'ultimo sistema non è causale perché per poter riportare l'output è necessario conoscere i valori dell'input al tempo $t+1$, che è nel
futuro.

\paragraph{Memoria:} Un sistema senza memoria calcola l'output al tempo $t$ usando esclusivamente il valore dell'input al tempo $t$. Inoltre, un sistema
causale con memoria non può richiederle di conoscere il futuro, come dice la proprietà \[
	y(t) = f(x(t))
\]

\paragraph{Stabilità:} Una proprietà che vale per l'input prima della trasformazione vale anche dopo per l'output. Esistono più tipi di stabilità, come
quella asintotica, dove graficamente il segnale converge verso lo $0$ andando avanti nel tempo, oppure quella bounded-input bounded-output (BIBO), la
quale richiede che il segnale si trovi sempre all'interno di un certo intervallo.

\paragraph{Invertibilità:} Un sistema è invertibile quando il suo output contiene informazioni sufficienti per ricostruire l'input.\newline

\noindent Esiste in particolare una certa classe di sistemi chiamata \textbf{lineari tempo-invarianti} (LTI), i quali detengono le omonime proprietà che
gli consentono di avere una rappresentazione mediante la convoluzione.\par
Per capire che cosa fa il sistema è necessario introdurre la \textbf{risposta all'impulso} $h(t)$, ovvero il suo output quando il segnale di entrata è
un delta di Dirac $\delta(t)$: \[
	h(t) = T\{\delta(t)\}
\]

\noindent Sappiamo che ogni input può essere rappresentato come una collezione pesata di impulsi traslati nel tempo: \[
	x(t) = \int_{-\infty}^{+\infty} x(\tau)\delta(t-\tau) d\tau = x(t)*\delta(t)
\]

\noindent E per i sistemi LTI, ogni impulso produce una copia di $h(t)$ e la loro linearità ci consente di sommare tutte le risposte ottenute. Ne
consegue che: \[
	y(t) = T{x}(t) = \int_{-\infty}^{+\infty} x(\tau)h(t-\tau) d\tau = (x*h)(t)
\]

\noindent Dunque un sistema LTI con risposta all'impulso $h(t)$ è sempre una convoluzione con dell'input con $h$, la quale agisce come \textbf{kernel}
del filtro ed è capace di determinare completamente il comportamento di I/O del sistema quando lo si considera in uno stato iniziale di riposo. La
formula corrispondente nel dominio delle frequenze è: \[
	y[n] = (x*h)[n] = \sum_{m = -\infty}^{\infty} x[m]h[n-m] =
		\sum_{m=-\infty}^{\infty} h[m]x[n-m]
\]

\noindent Dove $h[n]$ è la risposta all'impulso sul campione $n=0$, quindi centrata sull'origine. Se si considera la trasformata di Fourier di $h(t)$
otteniamo la \textbf{risposta in frequenza}, la quale non cambia internamente e fornisce una descrizione del sistema LTI nel dominio delle frequenze: \[
	H(\omega) = \mathcal{F}\{h(t)\} = \frac{Y(\omega)}{X(\omega)}
\]

\noindent Sappiamo che per il teorema della convoluzione, nel dominio delle frequenze si passa da convoluzione a prodotto, facendoci ottenere \[
	Y(\omega) = H(\omega)\cdot X(\omega)
\]

\noindent Ne concludiamo che i sistemi LTI sono descritti dalla convoluzione nel tempo e dal prodotto nelle frequenze; in particolare, $H(\omega)$,
detta anche \textbf{funzione di trasferimento}, fornisce informazioni in merito ai cambi di fase e la scala. Sarà il nostro punto di partenza per
definire i filtri.\par
Per un filtro LTI a tempo discreto, la risposta all'impulso $h[m]$ è il suo output quando il segnale in entrata è un impulso. Dunque per ogni input
$x[n]$, la risposta è: \[
	y[n] = \sum_{m = -\infty}^{\infty} h[m]x[n-m]
\]

\noindent A partire da ciò possiamo introdurre alcuni filtri notevoli che pesano le entrate con la funzione di trasferimento.

\paragraph{Filtro media mobile:} Esegue la media degli ultimi $M$ campioni. Generalmente applicato nel dominio del tempo, smussa il segnale attenuando
variazioni di frequenza elevata. Si definisce come: \[
	h[m] = \begin{cases}
		1/M & m = 0,1, ..., M-1\\
		0 & \text{Altrimenti}
	\end{cases}
\]

\noindent E l'output risultante è: \[
	y[n] = \frac{x[n] + x[n-1] + ... + x[n-M+1]}{M}
\]

\noindent Il filtro controlla sequenzialmente ogni campioni ed effettuerà la media tra quello corrente e tutti i precedenti. Per esempio, con $M=3$ si
prenderanno gli ultimi $3$ campioni e ne si calcolerà la media, con $M=4$ lo si farà con $4$ e così via. Come conseguenza diretta, maggiore sarà tale
valore, tanto più smussato sarà il segnale di output.\par
Notare infine che l'applicazione di questo filtro implica anche la perdita di alcune informazioni, poiché lo smussamento del segnale limita il range
delle frequenze considerata e alcune di loro saranno quindi rimosse.

\paragraph{Filtro gaussiano:} Similmente alla media mobile, calcola una media dei valori del segnale e dà più peso a quelli vicini al centro della
funzione. Nel dominio delle frequenze è definito come: \[
	h[m] = \frac{e^{-m^2/2\sigma^2}}{\sum_{r=-K}^K e^{-r^2/2\sigma^2}}
\]

\noindent Il denominatore ha lo scopo di normalizzare i pesi, sicché la loro somma risulti $1$ e che il peso sia sempre compreso tra $[0,1]$; l'output
risultante è: \[
	y[n] = \sum_{m = -K}^K h[m]x[n-m]
\]

\paragraph{Filtro di Notch:} Classe di filtri che rimuove una frequenza digitale specifica $\omega_0$. Un'impostazione comune definisce la propria
risposta in frequenza come: \[
	H(\omega) = \frac{1-2cos(\omega_0)z^{-1}+z^{-2}}{1-2rcos(\omega_0)z^{-1}+r^2z^{-2}};\quad \text{Con } 0 < r < 1
\]

\noindent Ciò che fa il filtro è creare degli zeri sulle frequenze da rimuovere; il denominatore controlla quanto stretto è l'intervallo di frequenze
$\omega_0$. Più vicino è a $1$, più stretto sarà.

\paragraph{Filtro passa-basso ideale:} Filtro che rimuove tutte le frequenze al di
sopra di un certo limite $f_c$. Una sua versione ideale ha una risposta in frequenza
definita come: \[
	H(f) = \begin{cases}
		1, |f| \leq f_c\\
		0, |f| > f_c
	\end{cases}
\]

\noindent La cui risposta all'impulso è una funzione sinc: \[
	h(t) = 2f_c sinc(2f_ct),\quad sinc(u) = \frac{sin(\pi u)}{\pi u}
\]

\noindent Questo filtro rappresenta un utile modello teorico, tuttavia si estende infinitamente in ambo le direzioni del tempo, rendendolo non causale
e impossibile da implementare esattamente in tempo reale; nella pratica viene infatti approssimato. Inoltre, esiste anche il complementare 
\textbf{passa-alto}, che rimuove tutte le frequenze al di sotto del taglio.

%

\section{Digitalizzazione dei segnali}
Per poter elaborare i segnali è necessario che questi siano convertiti da analogici a digitali tramite un convertitore \textbf{analog-to-digital} (ADC),
il quale segue i passi: \begin{enumerate}
	\item \textbf{Misurazione del segnale}: Un sensore trasforma la quantità fisica del segnale in un voltaggio $x(t)$ che può cambiare nel tempo.
	\item \textbf{Filtraggio per rimuovere frequenze troppo alte}: Prima del campionamento è necessario utilizzare un filtro passa-basso per rimuovere
        le frequenze che l'ADC non riesce a captare in modo affidabile. L'output è: \[
			y(t) = (x*h)(t)
	    \]
	\noindent Ovvero una convoluzione dell'input con la risposta all'impulso. Previene l'aliasing.
	\item \textbf{Scelta della frequenza di campionamento e misurazione}: L'ADC misura il voltaggio filtrato a intervalli regolari $T$, quindi a una
        frequenza di $F_s = 1/T$, facendo risultare la sequenza: \[
		    y[n] = y(nT)
	    \]
	\noindent Se la frequenza più alta da preservare è $B$, il minimo teorico assumibile dalla frequenza di campionamento per poter ricostruire il
        segnale fedelmente è $f_s > 2B$\footnote{Questa è la condizione di Nyquist nel teorema di Nyquist-Shannon, approfondita più avanti}.
	\item \textbf{Quantizzazione di ogni misura}: L'ADC arrotonda ogni voltaggio misurato a un determinato range di interi dato dal numero di bit
        assegnati. Con $b$ bit per campione ci sono $2^b$ codici possibili.
	\item \textbf{Encoding a bit}: Ogni valore quantizzato è salvato come numero binario, la cui sequenza rappresenta tutte le informazioni discusse
        prima.
\end{enumerate}

\noindent Introduciamo ora il \textbf{Teorema di Nyquist-Shannon}, il quale afferma che un segnale continuo a banda limitata può, in teoria, essere
perfettamente ricostruito da campioni uniformi se la frequenza di campionamento è più grande del doppio della sua frequenza più alta $f_{max}$\footnote{
Il valore della frequenza massima determina anche la banda $B$.}, ovvero: \[
	f_s > 2f_{max}
\]

\noindent Dove $f_s$ è la frequenza di campionamento. Campionare troppo lentamente causa il fenomeno dell'\textbf{aliasing}, dove due componenti
frequenziali si sovrascrivono l'una con l'altra, diventando indistinguibili. Inoltre: \begin{itemize}
	\item \textbf{Rateo di Nyquist} $2f_{max}$: Rateo minimo necessario per evitare aliasing quando un segnale la cui massima frequenza è $f_{max}$
        viene campionato.
	\item \textbf{Frequenza di Nyquist} $f_s/2$: Frequenza massima che il sistema di campionamento può rappresentare senza aliasing.
\end{itemize}

% TODO Riprendi dalla lezione 8